\scanpage{34}
\mspage{15}
\[
(53)\qquad a_i^\omega\varphi_i^\omega(g)=\varphi_i^{\omega'}(\varepsilon_0(g))
\]
Giving an endomorphism of the \inserted{profinite} groupoid, fixing the vertices and commuting with $G_J$, is equivalent to giving the \deleted{elements} images of the $a_i^\omega,c_i^\omega,d_i^\omega$, in a manner \deleted{which} compatible with the action of $G_J$, and in such a way as to satisfy the relations (48).

We shall have
\[
(54)\qquad\left\{\begin{aligned}
u(a_i^\omega)&=a_i^\omega\varphi_i^\omega(\lambda)\\
u(c_i^\omega)&=\varphi_{\omega'i}^{\omega'}(\eta)c_i^\omega=c_i^\omega\psi_{0}^\omega(\eta)\\
u(d_i^\omega)&=\psi_i^\omega(\zeta)d_i^\omega=d_i^\omega\widetilde\psi_i^\omega(\zeta)
\end{aligned}\right.
\]
with $\lambda,\zeta,\eta\in\widehat\Pi_{0,3}$, subject to the sole relation deduced from (48). Instead of redoing the calculation, somewhat tedious, to make explicit the relation in $(\lambda,\eta,\zeta)$, I prefer to note that the \deleted{the} groupoid under consideration is the union of the groupoids $\overline{\underline{\mathrm{II}}}$ and $\overline{\underline{\mathrm{IV}}}$, whose intersection is $\overline{\underline{\mathrm{III}}}$, hence \inserted{giving} an endomorphism \inserted{(fixing the vertices and commuting with $G_J$)} amounts to giving \inserted{such} an endomorphism $(\lambda,\eta)$ of $\overline{\underline{\mathrm{II}}}$, and another such endomorphism \inserted{$\zeta$} of $\overline{\underline{\mathrm{IV}}}$, in such a way that they glue together \deleted{\unsure{in} in $\overline{\underline{\mathrm{III}}}$} by a $=$ endomorphism $(\widetilde\alpha)$ of $\overline{\underline{\mathrm{III}}}$, i.e such that we have both
\[
(55)\qquad\widetilde\alpha=\eta^{-1}\varepsilon_0(\lambda\eta)\qquad\text{and}\qquad
\widetilde\alpha=\zeta\rho^{-1}(\zeta^{-1})
\]
\margin{33}
